\documentclass[11pt]{article}
\usepackage{amsmath,amssymb,hyperref}
\title{When clearing one claim forces another withdrawal}
\author{}
\date{}
\begin{document}
\maketitle

Withdrawing a claim from one handoff can expose a blocking claim at another. For a restricted model with fixed queues, every necessary withdrawal can be derived exactly. The process stops at the least set that clears all protected handoffs. Separate components can compute their own part, but they must exchange newly forced shared withdrawals until neither side forces more. Equal shared results do not imply equal private work or authority.

\paragraph{Operating rule.}
Each protected handoff sees a finite list of claims in fixed order. The first remaining claim decides whether that list blocks the handoff; an empty list does not block. A claim has a fixed bit at each occurrence: would it block if it became first? The same identity may occur at several protected handoffs, with a different bit at each. One actual withdrawal of that identity removes all its occurrences in the declared scope, once. This is permitted only for claims that the authorizing actor can withdraw and that have not already created a separate physical obligation. The supplied seed set of withdrawals is required. There are no additions, reorders, cell releases or changes of eligibility while applying this modeled batch. These assumptions strengthen the one-list-per-independent-claim model in \href{https://github.com/eternal-november/eos-research-programmes/pull/431}{the SUMO source comparison}; they do not assert that SUMO exposes or implements such batch withdrawals.

\paragraph{Exact propagation.}
For a blocking claim $x$ of a list, impose the implication: if every earlier claim in that list is removed, then $x$ must be removed. If an implication fails, $x$ is a surviving first blocker. Conversely, a surviving first blocker violates its implication. Thus clearing all the lists is equivalent to satisfying every implication and including the seed. The solutions are closed under intersection; the intersection of all solutions exists, since removing all claims is a solution before permission and budget constraints. It is the unique least withdrawal set. Starting at the seed, repeatedly withdrawing a current blocker is sound: every clearing solution already contains all withdrawals made so far, hence must also contain this one. When no blocker is first, the current set is itself a solution and so is least. This identifies exactly what creates the next dependency and what stops propagation.

For example, one handoff sees $a$ as a blocker. Another has $a$ as a nonblocker ahead of blocking $b$. Clearing the first forces withdrawing $a$, making $b$ unavoidable at the second. Protected $b$ proves no lawful joint clearing set. In the longer chain, every next local list starts with the previous claim as nonblocker, then the next claim as blocker. Any one of those holders may become the first protected claim, with every earlier withdrawal still necessary. A budget or permission failure can be decided exactly from the least set, since enlarging that set cannot repair an unauthorized withdrawal or lower a monotone cost. This applies to a shared withdrawal charged once. No inference about later physical effects follows unless the authorization premises actually cover every removed occurrence.

\paragraph{Assembly across components.}
Let $U_i$ contain all identities read by component $i$ and let $B$ include every shared identity. For current boundary withdrawals $X$, compute each component's same least closure with local seed and incoming $X\cap U_i$. Publish only withdrawals in $B$. Repeatedly union those with $X$ until this adds nothing; the finite round process is increasing. At each round every global clearing solution contains the boundary additions, by local leastness. At the fixed point, the local sets agree on shared identities: one side would otherwise have published a withdrawal not yet included on the other. Their union restricts to each local set and clears every list. Leastness follows from the preceding induction. A revision that changes queue order, effects or the occurrence scope requires rerunning or invalidating the affected component; an old summary is not a change detector.

If the other components can choose arbitrary boundary seeds and only observe the returned boundary set, the boundary closure function is an exact replacement interface. Equality lets the same iterations proceed with the same boundary results; a differing value is observable with a partner that seeds the distinguishing input and adds no new consequences. If a consumer also decides permission or allocates funds, the interface must additionally account for private violations and private cost at the stabilized input. Shared withdrawals are charged once. For the budget statement, costs are nonnegative additive across distinct identities and private identities are disjoint. As a counterexample, the queue with nonblocking boundary $a$ then blocking private $b$ exports the same boundary closure as the queue with nonblocking $a$ then nonblocking private $c$. Under seed $a$, the first additionally requires $b$; a zero private budget or forbidden $b$ distinguishes them. Boundary equality alone cannot prove reuse of the budget check.

\paragraph{Replay and comparisons.}
The head-cursor implementation watches each currently nonblocking live head and wakes that list when the head is withdrawn. Across a monotonically growing, fixed batch, each head is scanned at most twice, and each watching registration is woken at most once. Total local work is linear in the number of lists and inspected entries with unit-cost dictionary access, without charging discovery or changes to the queues. On the length-4,096 forced chain it inspects 12,286 entries, including the repeated head checks. This is not a bound on cost when implementations repeatedly recompute a whole interface snapshot, or when the shared graph itself changes.

The independent exact reference enumerates legal clearing sets and intersects them. The optimized direct Horn baseline shares a prefix atom between successive entries and applies standard forward chaining, also linear in the explicit queue input. It matches the head solution, so there is no gain over that direct representation. \href{https://core.ac.uk/download/pdf/82477946.pdf}{Dowling and Gallier} give linear Horn algorithms; no claim of novelty about Horn closure is made. The 1,600 finite families compare whole-set enumeration, direct Horn, local cursors and the boundary iteration, testing 34,149 seed supersets and 60,836 permission/budget alternatives. Three Lean results check the abstract restriction, least joining and budget rules; the queue-to-implication interpretation is proved above and exercised by replay. The operating result is a scoped propagation and replacement law, not an implementation of fault discovery, a live proof transfer under actual plant variation, or R2/I1 acceptance.
\end{document}
